% Auto-generated by tools/analyze_curvature_fg_full_validation.py
\newcommand{\MaxCurvatureChange}{0.00787753}
\newcommand{\CurvatureGateConclusion}{四个半径中有 4/4 个半径的六类中心位移全部通过 1\% 门槛，其中 4/4 个也通过 0.5\% 优先门槛；全表最大变化为 0.00787753\%。}
\newcommand{\FullRadiusStatus}{\pass}
\newcommand{\InteractionConclusion}{12 个半径--载荷交互中有 12 个高于三倍离散界；绝对值最大者位于 $R=1.0$ m 的电致工况，为 1.5421 个百分点。同一半径电致与磁致交互的最大差仅为 1.602e-08 个百分点，符合线性结构响应下的归一化一致性；机械工况有 4/4 个半径达到已分辨标准。}
\newcommand{\RadiusSummaryRows}{%
1.0 & 1.0 & 0.0079 & 0.0042 & 通过 \\
0.4 & 2.5 & 0.0059 & 0.0025 & 通过 \\
0.3 & 3.3 & 0.0053 & 0.0027 & 通过 \\
0.2 & 5.0 & 0.0046 & 0.0026 & 通过 \\
}
\newcommand{\InteractionRows}{%
1.0 & 机械 & -0.0999 & 0.0079 & 12.7 & 已分辨 \\
1.0 & 电致 & 1.5421 & 0.0032 & 479.6 & 已分辨 \\
1.0 & 磁致 & 1.5421 & 0.0032 & 479.6 & 已分辨 \\
0.4 & 机械 & -0.0473 & 0.0059 & 8.1 & 已分辨 \\
0.4 & 电致 & 1.4942 & 0.0015 & 977.1 & 已分辨 \\
0.4 & 磁致 & 1.4942 & 0.0015 & 977.1 & 已分辨 \\
0.3 & 机械 & -0.0362 & 0.0053 & 6.8 & 已分辨 \\
0.3 & 电致 & 1.4285 & 0.0020 & 706.4 & 已分辨 \\
0.3 & 磁致 & 1.4285 & 0.0020 & 706.4 & 已分辨 \\
0.2 & 机械 & -0.0266 & 0.0046 & 5.7 & 已分辨 \\
0.2 & 电致 & 1.3568 & 0.0023 & 596.9 & 已分辨 \\
0.2 & 磁致 & 1.3568 & 0.0023 & 596.9 & 已分辨 \\
}
